\chapter{Integrated Claim Ledger and Algebraic Research Program}

The previous chapters built the theory in layers. This chapter records how the
project artifacts in \texttt{temp/} and the companion drafts map into that
theory. Its purpose is traceability. A doctoral thesis cannot rely on the
statement that a formula was ``modified to taste.'' It must state where the
formula comes from, what model it assumes, what it proves, what it only
screens, and what validation gate is still missing.

The central integrated object is
\begin{equation}
  T(s;\rho)
  =
  s^2M(\rho)+sD_0(\rho)+L(\rho)+\Pi(s;\rho),
  \qquad
  \Pi(s;\rho)=\Pi_K(s;\rho)+s\Pi_D(s;\rho).
  \label{eq:integrated_operator}
\end{equation}
Here \(L\) is the operating-point stiffness inherited from power flow,
\(M\) is physical or virtual inertia, and \(\Pi(s)\) is the condensed
controller self-energy. The damping-channel specialization defines
\(\Sigma(s)=D_0+\Pi_D(s)\) when \(\Pi_K\) is absent or explicitly absorbed into
stiffness. The fixed-\(D\) theory is the special case \(\Pi(s)=0\) with direct
damping \(D_0=D\). The uniform proportional damping floor is the smaller
special case \(D=\gamma M\), while the broader decoupled fixed-\(D\) class is
\([\Lt,\Dt]=0\). This single inclusion chain is the clean way to organize the
project:
\begin{equation}
  \begin{aligned}
  D=\gamma M
  &\subset
  \{D:\ [\Lt,\Dt]=0\}\\
  &\subset
  \{D:\ [\Lt,\Dt]\neq0,\ D\ {\rm fixed}\}\\
  &\subset
  \{\Pi(s)\ {\rm rational}\}\\
  &\subset
  \{{\rm nonlinear\ memory}\}.
  \end{aligned}
  \label{eq:integrated_inclusion_chain}
\end{equation}

\thesisbox{thesisblue}{Chapter visual promise}{
This chapter turns the project archive into a controlled research program. The
reader should be able to trace each idea from source, to operator, to margin,
to action, and finally to a validation gate.}

\begin{table}[ht]
\centering
\caption{Claim-status vocabulary used in this chapter.}
\label{tab:claim_status_vocabulary}
\begin{tabular}{L{0.19\textwidth}L{0.67\textwidth}}
\toprule
\rowcolor{softblue}
\thead{Status} & \thead{Meaning}\\
\midrule
\statuspill{thesisgreen}{Algebraic result} &
The statement follows from matrix identities or perturbation formulas under
declared assumptions.\\
\statuspill{thesisblue}{Screen} &
The calculation ranks candidates or flags risk, but does not by itself certify
the protected margin.\\
\statuspill{thesisorange}{Gate} &
The result is an internal validation or rejection test that constrains what can
be claimed next.\\
\statuspill{thesisred}{Candidate} &
The mechanism is mathematically coherent but still requires a registered
benchmark before it becomes an empirical contribution.\\
\bottomrule
\end{tabular}
\end{table}

\FloatBarrier

\section{Source-to-Claim Ledger}

Tables~\ref{tab:source_claim_ledger_main} and
\ref{tab:source_claim_ledger_validation} convert the reviewed material into
thesis claims. The language is deliberately conservative: an algebraic
derivation is not a benchmark validation, and a validated internal gate is not
an industry-wide theorem.

\begin{table}[p]
\centering
\small
\caption{Ledger linking the main project sources to mathematical objects and
claim status.}
\label{tab:source_claim_ledger_main}
\begin{tabular}{L{0.17\textwidth}L{0.25\textwidth}L{0.30\textwidth}L{0.12\textwidth}}
\toprule
\rowcolor{softblue}
\thead{Source} & \thead{Mathematical object} & \thead{Incorporated thesis claim} & \thead{Status}\\
\midrule
Handoff protocol &
\(\zeta_k\simeq\delta_k/(2\sqrt{\nu_k})\), adaptive reduced QEP, line
sensitivity warnings &
Graph-frequency metrics are screens; damping-margin rankings must use
post-action \(\zeta_{\min}\) or pole recomputation. &
Calibrated protocol\\
\addlinespace
Consolidated blueprint &
Damping-region inequality
\(\delta_k\ge 2\zeta_\star\sqrt{\nu_k}\), reduced QEP selection, exact QEP
sensitivity &
Use the diagonal damping region when modes are separated, then promote only
coupled low-damping clusters to a reduced QEP. &
Theory plus surrogate gate\\
\addlinespace
Unified attribution note &
\(T(s)=s^2M+sD_0+L+\Pi(s)\), self-energy contribution, weak nodes and weak
links &
The final reduced controller object for IBR-rich models is rational in \(s\);
its damping-channel projection is \(\Sigma(s)\), not generally a constant
\(D\). &
Core thesis frame\\
\addlinespace
Hidden-resonance note &
Distance \(d^{\rm ctrl}_{kj}\), self-energy dominance \(\phi_\Sigma^{\rm abs}\), hidden
action margin &
A stable pole may be close to a controller-resonant damping loss even when its
present damping ratio is acceptable. &
Candidate diagnostic\\
\addlinespace
Line-shunt derivation &
\(\operatorname{Sym}(n)=\mathcal L_{\rm all}\oplus\mathcal D_{\rm sh}\) and
\(\partial\Lt/\partial M_i\) &
Virtual inertia has an existing-line plus shunt structural footprint in the
inertia-normalized stiffness; it does not create new-line directions. &
Algebraic result\\
\bottomrule
\end{tabular}
\end{table}

\begin{table}[p]
\centering
\small
\caption{Ledger linking internal drafts and validation records to incorporated
claims.}
\label{tab:source_claim_ledger_validation}
\begin{tabular}{L{0.17\textwidth}L{0.25\textwidth}L{0.30\textwidth}L{0.12\textwidth}}
\toprule
\rowcolor{softblue}
\thead{Source} & \thead{Mathematical object} & \thead{Incorporated thesis claim} & \thead{Status}\\
\midrule
Hosting-capacity draft PDF &
Mass-normalized Laplacian, first-order spectral hosting screen, stochastic fragility,
edge and node fragility &
Dynamic hosting capacity should be a set constrained by damping, resolvent,
and stochastic frequency-security margins, not only voltage limits. &
Internal draft, integrated cautiously\\
\addlinespace
Beyond-voltage draft PDF &
Kirchhoff-consistent AC power flow, weak nodes, hidden dynamic margins &
Voltage-safe operating points can still violate dynamic spectral margins; the
power-flow Jacobian is the bridge from electrical operating point to \(L\). &
Internal draft, integrated cautiously\\
\addlinespace
Frequency-estimator draft and ZIP &
\(T_{\rm risk}\), latency, peak-error exposure &
A dynamic hosting certificate should include what the protection and
measurement layer can actually observe. &
Protection gate\\
\addlinespace
Braess-NEP phase &
NEP line sensitivity with self-energy contribution &
Some line reinforcements are control-resonant damping risks, not generic
topological Braess theorems. &
Internal NEP gate\\
\addlinespace
Rational-filter output &
Diagonal, second-order, reduced-QEP, and adaptive errors &
The adaptive estimator is useful because diagonal screens are often accurate
but fail under clustered or strongly coupled damping. &
Surrogate validation\\
\addlinespace
Coordination-lever gate &
Equal-cost topology, damping, and inertia packages &
Mixed lever packages can outperform the best single lever, but modal and
resolvent margins need separate reporting. &
Planning gate\\
\bottomrule
\end{tabular}
\end{table}

\FloatBarrier

\section{Step-by-Step Analysis Stack}

The analysis should be read as a sequence. Each step adds one layer of physics
or mathematics.

\begin{enumerate}
\item \textbf{Start from the electrical operating point.} Bus voltages,
current injections, and admittances satisfy \(I=Y_{\rm bus}V\). Complex power
is \(S_i=P_i+jQ_i=V_iI_i^*\). This gives the nonlinear AC power-flow map.
\item \textbf{Linearize active power with respect to angle.} At a feasible
operating point, the block \(J_{P\theta}=\partial P/\partial\theta\) gives the
synchronizing stiffness. After the reference-angle direction is removed, write
the positive stiffness as \(L\). In the lossless small-angle case, \(L\) is a
weighted Laplacian.
\item \textbf{Attach nodal dynamics.} Structure-preserving power-system models
retain bus-level dynamics instead of hiding all loads inside a reduced network
\cite{bergen1981structure,kundur1994power,sauer2017power,milano2010power}.
For the electromechanical layer this gives
\[
  M\ddot\theta+D\dot\theta+L\theta=p.
\]
\item \textbf{Remove the uniform angle mode.} Work on a subspace orthogonal to
the rotational symmetry. If \(R^T\mathbf 1=0\), define
\[
  M_r=R^TMR,\qquad D_r=R^TDR,\qquad L_r=R^TLR.
\]
\item \textbf{Normalize inertia.} With
\[
  \Lt=M_r^{-1/2}L_rM_r^{-1/2},\qquad
  \Dt=M_r^{-1/2}D_rM_r^{-1/2},
\]
the QEP becomes
\[
  \left(s^2I+s\Dt+\Lt\right)v=0.
\]
\item \textbf{Diagonalize stiffness, not damping.} Let
\(\Lt Q=Q\operatorname{diag}(\nu_k)\). The modal damping matrix is
\[
  \Gamma=Q^T\Dt Q.
\]
Only when \(\Gamma\) is diagonal do graph modes decouple exactly.
\item \textbf{Use the diagonal damping screen first.} The first screen is
\[
  \hat\zeta_k=\frac{\Gamma_{kk}}{2\sqrt{\nu_k}},
  \qquad
  h_k=\Gamma_{kk}-2\zeta_\star\sqrt{\nu_k}.
\]
It is exact under proportional or commuting damping, and a first-order
diagnostic otherwise.
\item \textbf{Promote dangerous clusters to a reduced QEP.} If off-diagonal
modal damping is large or the graph frequencies are clustered, solve the QEP
on the selected modal subspace.
\item \textbf{Condense controllers only after admitting that the object
changes.} Eliminating controller states produces \(\Pi(s)\), with
\(\Sigma(s)\) only as a declared damping-channel specialization, so the problem
becomes the rational NEP \eqref{eq:integrated_operator}.
\item \textbf{Plan with margins, not labels.} A bus, line, or controller is
weak only relative to a margin and an action family: \(\zeta_{\min}\), protected
resolvent gain, control-pole distance, stochastic variance, or protection
exposure.
\end{enumerate}

\begin{figure}[ht]
\centering
\begin{tikzpicture}[
  nodebox/.style={
    rectangle,
    rounded corners=3pt,
    draw=#1!70!black,
    fill=#1!10,
    text width=0.22\textwidth,
    minimum height=1.05cm,
    align=center,
    font=\small
  },
  arr/.style={-{Latex[length=2.2mm]}, thick, draw=thesisgray}
]
\node[nodebox=thesisblue] (op) {\textbf{Operating point}\\
\(V^\star,\theta^\star\)};
\node[nodebox=thesisgreen, right=0.45cm of op] (lm) {\textbf{Graph stiffness}\\
\(L,\Lt,\nu_k,q_k\)};
\node[nodebox=thesisorange, right=0.45cm of lm] (dm) {\textbf{Damping screen}\\
\(\delta_k,\hat\zeta_k,h_k\)};
\node[nodebox=thesisred, below=0.7cm of lm] (nep) {\textbf{Controller bridge}\\
\(\Sigma(s),s\Sigma'(s)\)};
\node[nodebox=thesisgold, right=0.45cm of nep] (act) {\textbf{Action ranking}\\
\(\Delta w,\Delta M,\Delta D,\Delta\kappa\)};
\draw[arr] (op) -- (lm);
\draw[arr] (lm) -- (dm);
\draw[arr] (dm.south) -- ++(0,-0.38) -| ([xshift=0.24cm]act.north);
\draw[arr] (lm.south) -- (nep.north);
\draw[arr] (nep) -- (act);
\end{tikzpicture}
\caption{Analysis stack used by the integrated research program. The diagonal
damping screen is the fast route; the controller bridge is invoked when static
damping cannot preserve the pole-motion mechanism.}
\label{fig:integrated_analysis_stack}
\end{figure}

\FloatBarrier

\section{Adaptive Reduced-QEP Algebra}

The off-diagonal entries of \(\Gamma\) are the algebraic reason the scalar
screen can fail. A compact coupling diagnostic is
\begin{equation}
  \eta_{\rm off}
  =
  \frac{\norm{\Gamma-\operatorname{diag}(\Gamma)}_F}{\norm{\Gamma}_F},
  \label{eq:eta_off}
\end{equation}
and a mode-pair diagnostic is
\begin{equation}
  c_{k\ell}
  =
  \frac{|\Gamma_{k\ell}|}
       {|\sqrt{\nu_k}-\sqrt{\nu_\ell}|+\varepsilon_{\rm gap}}.
  \label{eq:pair_coupling}
\end{equation}
Large \(c_{k\ell}\) means that the damping coupling is not small relative to
the frequency separation. This is a direct application of matrix perturbation
logic \cite{stewart1990matrix,kato1995perturbation}: small denominators make
apparently small off-diagonal terms important.

The adaptive estimator used by the project is therefore:
\begin{enumerate}
\item Compute \((\nu_k,q_k)\) of \(\Lt\) and \(\Gamma=Q^T\Dt Q\).
\item Rank all modes by \(\hat\zeta_k=\Gamma_{kk}/(2\sqrt{\nu_k})\).
\item Select low-\(\hat\zeta_k\) modes and augment the set with neighbors whose
\(c_{k\ell}\) exceeds a declared threshold.
\item Build \(V=[q_k]_{k\in\mathcal I}\), \(L_{\mathcal I}=V^T\Lt V\), and
\(D_{\mathcal I}=V^T\Dt V\).
\item Solve
\[
  \det(s^2I_{\mathcal I}+sD_{\mathcal I}+L_{\mathcal I})=0.
\]
\item Report the reduced-QEP damping margin and the gap to the diagonal
screen. If the gap is large, the graph-only explanation is not enough.
\end{enumerate}

The \texttt{ieee39\_rational\_filter} surrogate record supports this workflow:
over 300 trials the diagonal screen had median relative error
\(4.68\times10^{-3}\), while the second-order correction and reduced-QEP
variants reduced the median to approximately \(3.2\times10^{-4}\) and
\(3.7\times10^{-4}\), respectively. This does not prove universal accuracy. It
does justify using the diagonal formula as a fast screen and the reduced QEP as
a targeted correction.

\section{Exact Lever Sensitivities}

Let \(p\) be any scalar action: line weight, local inertia, local damping,
converter gain, or hosting-penetration coordinate. For the fixed-\(D\) QEP
\[
  T_Q(s,p)=s^2I+s\Dt(p)+\Lt(p),
\]
with right and left eigenvectors \(v\) and \(y\), a simple pole satisfies
\begin{equation}
  \frac{ds}{dp}
  =
  -
  \frac{y^*\left(s\,\Dt_p+\Lt_p\right)v}
       {y^*\left(2sI+\Dt\right)v}.
  \label{eq:integrated_qep_sensitivity}
\end{equation}
This is the QEP version of the nonlinear eigenvalue sensitivity used for
matrix polynomials and NEPs \cite{tisseur2001qep,guettel2017nep,beyn2012integral}.
If \(s=\sigma+j\omega\), \(ds/dp=a+jb\), and
\(\zeta=-\sigma/|s|\), then
\begin{equation}
  \frac{d\zeta}{dp}
  =
  \frac{\sigma\omega b-\omega^2a}{|s|^3}.
  \label{eq:zeta_from_pole_sensitivity}
\end{equation}
This equation is important because a line can move a pole left in real part
and still reduce damping ratio if it increases oscillation frequency too much.

For a line \(e=(i,j)\) in the inertia-normalized graph layer,
\begin{equation}
  \frac{\partial\nu_k}{\partial w_e}
  =
  \left(
  \frac{q_k(i)}{\sqrt{M_i}}-
  \frac{q_k(j)}{\sqrt{M_j}}
  \right)^2.
  \label{eq:integrated_line_nu}
\end{equation}
But damping requires more than \(\partial\nu_k/\partial w_e\). With fixed
modal labeling,
\begin{align}
  \frac{\partial \zeta_k}{\partial w_e}
  &=
  \frac{1}{2\sqrt{\nu_k}}\frac{\partial\delta_k}{\partial w_e}
  -
  \frac{\delta_k}{4\nu_k^{3/2}}
  \frac{\partial\nu_k}{\partial w_e},
  \label{eq:integrated_line_zeta}\\
  \frac{\partial\delta_k}{\partial w_e}
  &=
  q_k^T\Dt_{w_e}q_k
  +
  2\sum_{\ell\ne k}
  \frac{
    (q_\ell^T\Lt_{w_e}q_k)(q_\ell^T\Dt q_k)
  }{\nu_k-\nu_\ell}.
  \label{eq:integrated_delta_rotation}
\end{align}
The first term in \eqref{eq:integrated_delta_rotation} is direct damping
change. The sum is modal rotation. It is zero under proportional damping and
can be decisive under heterogeneous damping.

For the controller-condensed NEP \eqref{eq:integrated_operator},
\begin{equation}
  \frac{ds}{dp}
  =
  -
  \frac{
  y^*\left(s^2M_p+sD_{0,p}+L_p+\Pi_p(s)\right)x
  }{
  y^*\left(2sM+D_0+\Pi_s(s)\right)x
  }.
  \label{eq:integrated_nep_sensitivity}
\end{equation}
Here the subscript \(p\) denotes a total derivative after re-solving the
operating point and updating the controller blocks. The additional term
\(\Pi_s(s)\), which reduces to \(\Sigma(s)+s\Sigma'(s)-D_0\) in the
damping-channel case, is the spectral memory of the condensed controller block.
It is exactly the term that a static damping bridge removes.

\subsection{Mechanism-Resolved Attribution}

Equation~\eqref{eq:integrated_nep_sensitivity} is not only a way to compute a
pole derivative. It is also an attribution identity once the total derivative
is logged by mechanism. For a declared operating-point continuation, write
\begin{equation}
  T_p
  =
  T_p^{\rm net}
  +T_p^{\rm op}
  +T_p^{\rm ports}
  +T_p^{\rm ctrl}
  +T_p^{M,D}.
  \label{eq:mechanism_tp_split}
\end{equation}
The terms mean, respectively: direct network change, equilibrium or power-flow
shift, change in controller port visibility and residues, movement of
condensed controller poles, and direct mass-damping change. This split is not
unique until the model assembly is fixed, but after that bookkeeping choice it
is algebraic.

Let \(s=\sigma+j\omega\) and define the real linear map
\begin{equation}
  \mathcal D_\zeta(s)[\eta]
  =
  \frac{\sigma\omega\,\Imag \eta-\omega^2\,\Real \eta}{|s|^3}.
  \label{eq:zeta_linear_map}
\end{equation}
For each channel \(c\in\{\rm net,op,ports,ctrl,MD\}\), set
\begin{equation}
  g_c
  =
  \mathcal D_\zeta(s)
  \left[
  \frac{-y^*T_p^c x}{y^*T_sx}
  \right].
  \label{eq:channel_zeta_gain}
\end{equation}
Then
\begin{equation}
  \frac{d\zeta}{dp}
  =
  g_{\rm net}
  +g_{\rm op}
  +g_{\rm ports}
  +g_{\rm ctrl}
  +g_{M,D}.
  \label{eq:zeta_channel_decomposition}
\end{equation}
The denominator \(y^*T_sx\) is shared by all channels. Therefore the rational
memory term \(\Pi_s(s)\) does more than add one extra numerator contribution:
it can rotate or amplify every channel simultaneously. This is the causal
language needed to turn a pole plot into a mechanism statement.

\begin{definition}[Strict controller-induced reinforcement reversal]
For a line or parameter action \(p\), a strict controller-induced reinforcement
reversal occurs when a network-strength score improves and the static
reduction does not predict damping degradation, but the controller-aware NEP
does:
\begin{equation}
  \frac{dG_{\rm graph}}{dp}>0,
  \qquad
  \frac{d\zeta_{\rm static}}{dp}\ge0,
  \qquad
  \frac{d\zeta_{\rm NEP}}{dp}<0.
  \label{eq:strict_controller_reversal}
\end{equation}
Equivalently, the static and controller-aware models disagree in sign:
\begin{equation}
  \operatorname{sign}
  \left(\frac{d\zeta_{\rm static}}{dp}\right)
  \ne
  \operatorname{sign}
  \left(\frac{d\zeta_{\rm NEP}}{dp}\right).
  \label{eq:static_nep_sign_mismatch}
\end{equation}
\end{definition}

This definition is intentionally stronger than the statement
\(d\lambda_2/dp>0\) and \(d\zeta/dp<0\). The latter can occur in an elementary
QEP because increasing stiffness may raise modal frequency faster than decay.
The stronger claim is that the usual static reduction can recommend the wrong
line or parameter action.

\begin{proposition}[Controller-induced reversal certificate]
Assume a simple protected pole and the decomposition
\eqref{eq:mechanism_tp_split}. Suppose the rational self-energy admits a local
pole-residue expansion
\begin{equation}
  \Pi(s,p)
  =
  \sum_\ell \frac{Q_\ell(s,p)}{s-\alpha_\ell(p)}
  +\Pi_{\rm rem}(s,p),
  \label{eq:pi_pole_residue_expansion}
\end{equation}
and the contribution of all neglected or grouped terms to \(d\zeta/dp\) is
bounded above by \(\bar r\). Let
\begin{equation}
  g_0=g_{\rm net}+g_{\rm op}+g_{M,D}
  \label{eq:base_static_gain}
\end{equation}
and let \(g_\ell\) be the damping-ratio contribution of one dominant
controller pole \(\alpha_\ell\):
\begin{equation}
  g_\ell
  =
  \mathcal D_\zeta(s)
  \left[
  \frac{-y^*\Pi_{p,\ell}(s)x}{y^*T_sx}
  \right],
  \qquad
  \Pi_{p,\ell}(s)
  =
  \frac{Q_{\ell,p}(s)}{s-\alpha_\ell}
  +
  \frac{Q_\ell(s)\alpha_{\ell,p}}{(s-\alpha_\ell)^2}.
  \label{eq:dominant_pole_reversal_gain}
\end{equation}
If
\begin{equation}
  g_0\ge0,
  \qquad
  g_0+g_\ell+\bar r<0,
  \label{eq:controller_induced_reversal_certificate}
\end{equation}
then the first-order degradation is certified as controller-induced rather
than merely controller-correlated.
\end{proposition}

\begin{proof}
Linearity of \eqref{eq:integrated_nep_sensitivity} and
\eqref{eq:zeta_linear_map} gives the channel sum
\eqref{eq:zeta_channel_decomposition}. The pole-residue derivative in
\eqref{eq:dominant_pole_reversal_gain} separates residue change from movement
of the condensed pole. If the non-controller baseline is nonnegative and the
dominant controller-pole term plus the worst admissible remainder is negative,
then every compatible completion of the grouped remainder has negative total
first-order damping-ratio derivative. The sign reversal is therefore forced by
the rational controller channel.
\end{proof}

\section{Hidden Control-Resonance Margin}

Let \(\mu_j\in\operatorname{spec}(A_{cc})\) be a condensed controller pole and
let \(s_k\) be a retained network-family pole. Define
\begin{equation}
  d^{\rm ctrl}_{kj}
  =
  \frac{|s_k-\mu_j|}{|s_k|+|\mu_j|},
  \qquad
  \phi_\Sigma^{\rm abs}(k)
  =
  \frac{|y_k^*s_k\Sigma'(s_k)x_k|}
  {|y_k^*2s_kMx_k|+
   |y_k^*\Sigma(s_k)x_k|+
   |y_k^*s_k\Sigma'(s_k)x_k|+
   \varepsilon_{\rm reg}}.
  \label{eq:hidden_ctrl_metrics}
\end{equation}
Small \(d^{\rm ctrl}_{kj}\) and large \(\phi_\Sigma^{\rm abs}(k)\) mean that a presently
stable pole lies in the resolvent shadow of a controller pole. If \(A_{cc}\) is
diagonalizable with eigenvector condition factor \(\kappa_c\), then outside
\(\operatorname{spec}(A_{cc})\)
\begin{equation}
  \norm{(sI-A_{cc})^{-1}}
  \le
  \frac{\kappa_c}{\Delta_c(s)},
  \qquad
  \Delta_c(s)=\min_j |s-\mu_j|,
  \label{eq:resolvent_bound}
\end{equation}
and
\begin{equation}
  \norm{\Sigma'(s)}
  \le
  \frac{\norm{C}\norm{B}\kappa_c^2}{\Delta_c(s)^2}.
  \label{eq:sigma_prime_bound}
\end{equation}
Thus the warning is algebraic, not rhetorical: as a retained pole approaches a
condensed controller pole, the self-energy slope can become large.
\(\phi_\Sigma^{\rm abs}\) is a dominance index, not a strict fraction; the
absolute denominator avoids artificial inflation or suppression caused by
complex cancellation in \(y^*T_sx\).

For an admissible action \(p\) and target margin \(\zeta_\star\), define
\begin{equation}
  m_p^\zeta(k)
  =
  \frac{\zeta_k-\zeta_\star}
       {\left[-\partial\zeta_k/\partial p\right]_+},
  \qquad
  m_p^d(k,j)
  =
  \frac{d^{\rm ctrl}_{kj}-d_\star}
       {\left[-\partial d^{\rm ctrl}_{kj}/\partial p\right]_+}.
  \label{eq:hidden_action_margins}
\end{equation}
The hidden margin is
\begin{equation}
  m_p^{\rm hid}(k,j)=\min\{m_p^\zeta(k),m_p^d(k,j)\}.
  \label{eq:hidden_margin}
\end{equation}
This quantity says how much of action \(p\) can be applied before either the
damping threshold or the controller-distance threshold is locally exhausted.
The \texttt{phase\_braess\_nep} record gives one internal mechanism gate: 460
line-mode perturbations produced 40 strict resonant cases, and the clearest
line action moved a network-family mode toward a PLL-dominated condensed pole
with a self-energy-dominated damping decrease. The correct thesis claim is
therefore control-resonant reinforcement risk, not a universal Braess paradox.
The next extraction must count the smaller but stronger subset satisfying
\eqref{eq:strict_controller_reversal}; that subset is the empirical target for
the controller-induced reversal paper.

\subsection{Static Adequacy Gate and Resolvent Shadow}

The practical question is when the fixed-\(D\) or static-QEP screen is
sufficient. Let \(\Omega_k\) be a disk or contour neighborhood of a protected
pole \(s_k\) with radius \(r_k\). A first-order staticization error indicator is
\begin{equation}
  \varepsilon_{{\rm static},k}
  =
  \kappa_k r_k
  \sup_{\xi\in\Omega_k}\norm{\Pi_s(\xi)}_2,
  \qquad
  \kappa_k
  =
  \frac{\norm{x_k}_2\norm{y_k}_2}
       {|y_k^*T_s(s_k)x_k|+\varepsilon_{\rm reg}}.
  \label{eq:static_adequacy_gate}
\end{equation}
If \(\varepsilon_{{\rm static},k}<\tau m_k^{\rm margin}\) for a declared
tolerance \(\tau\), then the static screen is acceptable for that mode-action
pair. Otherwise the candidate lies in the resolvent shadow and must be checked
with the Schur NEP or the full DAE model. This turns the hierarchy into a
computational policy: use the cheap QEP broadly, escalate only shadowed modes,
and reserve full-model recomputation for finalists.

Distance to a condensed pole is not enough when \(A_{cc}\) is non-normal. The
screen should combine port-visible resolvent gain, modal projection, rational
memory, and pole conditioning:
\begin{align}
  \chi_c(s_k)
  &=
  \norm{C(s_kI-A_{cc})^{-1}(B_\theta+s_kB_\omega)}_2,
  \label{eq:port_visible_resolvent_gain}\\
  \chi_{k\ell}^{\rm proj}
  &=
  \frac{|y_k^*Q_\ell(s_k)x_k|}
       {|s_k-\alpha_\ell|\,\norm{y_k}_2\norm{x_k}_2+\varepsilon_{\rm reg}},
  \label{eq:projected_resolvent_shadow}\\
  \rho_{{\rm mem},k}
  &=
  \frac{|y_k^*\Pi_s(s_k)x_k|}
       {|y_k^*(2s_kM+D_0)x_k|+|y_k^*\Pi_s(s_k)x_k|
        +\varepsilon_{\rm reg}}.
  \label{eq:memory_dominance_index}
\end{align}
The desired empirical result is not that every large value produces a
reversal. It is that strict sign reversals concentrate where
\(\chi_{k\ell}^{\rm proj}\), \(\rho_{{\rm mem},k}\), and \(\kappa_k\) are large.
That statement can be tested with precision, recall, false-negative rate, and
the share of strict reversals captured after reviewing only the top decile of
line actions by resolvent-shadow score.

\section{Stochastic Fragility and Protection Exposure}

The hosting draft also contains a stochastic interpretation. In first-order
state form
\begin{equation}
  \dot x=Ax+B\xi,\qquad z=Cx,
  \label{eq:stochastic_state}
\end{equation}
with white input \(\xi\), the stationary covariance \(X\) satisfies
\begin{equation}
  AX+XA^T+BB^T=0,
  \qquad
  \mathbb E\norm{z}^2=\operatorname{trace}(CXC^T).
  \label{eq:lyapunov_variance}
\end{equation}
This is the same Gramian object used in network coherence analysis
\cite{bamieh2012coherence,tyloo2018robustness}. In a homogeneous proportional
limit, the graph contribution reduces to Kirchhoff-index-like quantities such
as \(\operatorname{trace}(L^+)\), which explains why effective resistance can
describe average vulnerability. In heterogeneous damping or rational
self-energy models, however, the covariance must be computed from the actual
linearization or an explicitly justified reduction.

The protection layer adds a separate constraint. Let \(e[m]\) be a frequency
estimator error and \(e_{\rm th}\) a protection threshold. The estimator draft
uses
\begin{equation}
  T_{\rm risk}
  =
  \sum_{m=0}^{N-1}\Delta t\,
  \mathbf 1_{\{|e[m]|>e_{\rm th}\}},
  \label{eq:integrated_trip_risk}
\end{equation}
which measures how long a protection-relevant error is exceeded. A complete
dynamic hosting set can therefore be written as
\begin{align}
  \mathcal H_{\rm dyn}^{\rm full}
  =
  \{\,\rho\in\mathcal H_{\rm stat}:&
  \ \zeta_{\min}(\rho)\ge\zeta_\star,\ 
  g_{\rm prot}(\rho)\le g_\star,\nonumber\\
  &
  \operatorname{trace}(C X(\rho) C^T)\le \sigma_\star^2,\ 
  T_{\rm risk}(\rho)\le T_\star,\ 
  d^{\rm ctrl}_{kj}(\rho)\ge d_\star
  \,\}.
  \label{eq:full_dynamic_hosting_set}
\end{align}
The point is not that every study must include every constraint. The point is
that each omitted constraint must be named.

\section{Pole-Motion Cone and Planning Optimization}

Collect line, inertia, damping, and controller actions in
\begin{equation}
  u=
  \begin{bmatrix}
    \Delta w & \Delta M & \Delta D & \Delta\kappa
  \end{bmatrix}^T.
  \label{eq:integrated_action_vector}
\end{equation}
For each monitored pole \(s_k\), NEP sensitivity gives
\begin{equation}
  \Delta s_k\simeq J_k u,
  \qquad
  J_k=
  \begin{bmatrix}
    \partial s_k/\partial p_1 & \cdots & \partial s_k/\partial p_r
  \end{bmatrix}.
  \label{eq:pole_motion_cone}
\end{equation}
Through \eqref{eq:zeta_from_pole_sensitivity}, this induces
\[
  \Delta\zeta_k\simeq g_k^Tu.
\]
Feasible actions therefore form a cone of pole motions in the complex plane.
Safe actions lie in the damping-improving half-space. Hidden-resonant actions
may still improve a graph metric while moving a pole toward a dangerous
controller resolvent region.

\begin{figure}[ht]
\centering
\begin{tikzpicture}[
  scale=0.95,
  axis/.style={->, thick, draw=thesisgray},
  safe/.style={fill=softgreen, draw=thesisgreen!70!black, line width=0.45pt},
  risk/.style={fill=softred, draw=thesisred!70!black, line width=0.45pt},
  arrow/.style={-{Latex[length=2.2mm]}, thick},
  lab/.style={font=\small}
]
\draw[axis] (-4.2,0) -- (2.2,0) node[right,lab] {\(\Real(s)\)};
\draw[axis] (0,-2.6) -- (0,2.6) node[above,lab] {\(\Imag(s)\)};
\draw[dashed, thesisgray] (-0.95,-2.3) -- (-0.95,2.3)
  node[above,lab] {target damping boundary};
\path[safe] (-2.6,0.35) -- (-3.65,1.45) -- (-3.65,-0.75) -- cycle;
\path[risk] (-2.6,0.35) -- (-1.05,1.15) -- (-0.95,-0.25) -- cycle;
\fill[thesisblue] (-2.6,0.35) circle (2.1pt)
  node[above left,lab] {current pole \(s_k\)};
\fill[thesisred] (-1.2,1.5) circle (2.1pt)
  node[above right,lab] {controller pole \(\mu_j\)};
\draw[arrow, thesisgreen] (-2.6,0.35) -- (-3.25,0.95)
  node[midway,above left,lab] {safe action};
\draw[arrow, thesisred] (-2.6,0.35) -- (-1.45,1.05)
  node[midway,below right,lab] {hidden-resonant action};
\node[lab, align=center, text width=0.26\textwidth] at (-3.35,-1.85)
  {larger damping ratio\\and larger controller distance};
\node[lab, align=center, text width=0.26\textwidth] at (1.0,-1.65)
  {rightward motion or\\controller approach is rejected};
\end{tikzpicture}
\caption{Pole-motion cone induced by the sensitivity matrix \(J_k\). A graph
action is acceptable only if its pole-motion vector improves the protected
margin and does not move the mode into a controller-resonant region.}
\label{fig:pole_motion_cone_visual}
\end{figure}

A local planning problem follows:
\begin{align}
  \max_{u,t}\quad & t \label{eq:planning_problem}\\
  \text{subject to}\quad
  & \zeta_k+g_k^Tu\ge t,\qquad k\in\mathcal K, \nonumber\\
  & d^{\rm ctrl}_{kj}+h_{kj}^Tu\ge d_\star,\qquad (k,j)\in\mathcal R,
    \nonumber\\
  & c^T|u|\le B,\qquad u_{\min}\le u\le u_{\max}, \nonumber\\
  & \text{AC, thermal, voltage, and protection feasibility gates hold.}
    \nonumber
\end{align}
This problem states exactly what the project can do with the new analysis:
rank SG-to-IBR conversions by dynamic headroom, identify where virtual damping
or inertia is honest, find line upgrades that improve damping rather than only
stiffness, detect control-resonant line risks, and compare mixed packages of
topology, damping, inertia, and control under equal cost.

\section{What This Analysis Enables}

The framework is useful because it converts broad planning questions into
declared algebraic tasks. Table~\ref{tab:analysis_enables} summarizes the
main actions.

\begin{table}[ht]
\centering
\small
\caption{Planning and research tasks enabled by the integrated analysis.}
\label{tab:analysis_enables}
\begin{tabular}{L{0.22\textwidth}L{0.25\textwidth}L{0.25\textwidth}L{0.16\textwidth}}
\toprule
\rowcolor{softblue}
\thead{Task} & \thead{Algebraic object} & \thead{Output} & \thead{Required gate}\\
\midrule
SG-to-IBR hosting direction &
\(h_k(\rho)=\delta_k-2\zeta_\star\sqrt{\nu_k}\) &
Maximum conversion direction before damping headroom is exhausted &
Full QEP or NEP recomputation\\
\addlinespace
Virtual inertia placement &
\(\partial h_k/\partial M_i\), RoCoF and damping weights &
Bus ranking that separates frequency support from damping support &
Post-action \(\zeta_{\min}\) and time-domain response\\
\addlinespace
Line reinforcement without batteries &
\(\partial\nu_k/\partial w_e\), \(\partial\zeta_k/\partial w_e\) &
Line list that improves margin, not only graph stiffness &
Flow, voltage, and damping checks\\
\addlinespace
Control-resonant risk audit &
\(d^{\rm ctrl}_{kj}\), \(\phi_\Sigma^{\rm abs}(k)\), \(s\Sigma'(s)\) &
Flag for stable poles that are close to controller-mediated damping loss &
Schur NEP or full DAE eigenanalysis\\
\addlinespace
Mixed lever package design &
Pole-motion cone \(J_ku\) and budget \(c^T|u|\le B\) &
Equal-cost packages of topology, damping, inertia, and control &
Registered baseline comparison\\
\addlinespace
Protection-aware hosting &
\(T_{\rm risk}\), estimator latency, protected gain &
Hosting set that includes what protection can observe &
Estimator and relay exposure test\\
\addlinespace
Self-energy passivity retuning &
\(\mu_\Sigma(\Omega)=\inf_{\omega\in\Omega}
\lambda_{\min}H_\Sigma(\omega)\) and \(ds_\star/d\rho\) &
Controller-gain update that improves both a passivity screen and the protected
pole-motion direction &
Schur NEP plus full-model retuning gate\\
\addlinespace
Sheaf weak-element research &
Cellular-sheaf Laplacian with bus/controller stalks and line restriction maps &
Candidate weak-node or weak-link certificate that is richer than scalar
Laplacian sensitivity &
New benchmark showing information not recovered by \(L\), QEP, or
\(\Sigma(s)\) alone\\
\bottomrule
\end{tabular}
\end{table}

\FloatBarrier

The last two rows are deliberately not ordinary planning claims. The passivity
row is a near-term extension because it uses the existing Schur self-energy
operator and adds a controller-respected dissipativity screen. The sheaf row is
a longer-term research program. It may become a powerful structural language
only if it produces a weak-element certificate that scalar graph theory and the
present self-energy diagnostics cannot already reproduce.

\section{Three Go--No-Go Studies}

The next work should be staged sequentially. First, the sign-characteristic
bridge must be grounded in the exact hypotheses of the matrix-polynomial
literature. Second, the passivity-retuning result should be made operational
because it extends the existing \(\Sigma(s)\) framework and can fail cleanly.
Third, the cellular-sheaf idea should be tested as a falsifiable structural
extension, not adopted as a new language by default.

\subsection{Sign-Characteristic Gate}

The strongest possible theoretical headline is not that controller condensation
makes a rational operator. That is Schur-complement mathematics. The stronger
claim is that a planning action can move a network-family pole through the
resolvent of a condensed controller and produce a sign collision that a
constant-\(D\) model cannot predict.

This claim needs a gate before it becomes a theorem. The gate has four steps:
\begin{enumerate}
\item Read the sign-characteristic definitions in the Hermitian
matrix-polynomial and self-adjoint meromorphic literature and list their
assumptions: regularity, symmetry, leading coefficient, elementary divisors, and
structure-preserving transformations.
\item State exactly which assumptions fail for
\(T(s)=s^2M+sD_0+L+\Pi(s)\): rationality, dissipation, non-normality, and loss
of a fixed symplectic or Hermitian product.
\item Build a two-bus SCR benchmark with a synchronous source and a
grid-following converter. The reduced coefficients must come from the exact
linearized algebraic-differential Jacobian; placeholders such as ``or similar''
are not allowed.
\item Track three quantities while varying SCR or PLL bandwidth:
\(\chi_k^{\rm op}\), \(\phi_\Sigma^{\rm abs}\), and the full-model eigenvalues.
The gate passes only if the same event appears as signature change or
singularity, self-energy dominance, and damping loss.
\end{enumerate}

If this gate passes, the first Transactions-style paper should be narrowly
framed around controller-resolvent reinforcement reversals:
\[
  \frac{dG_{\rm graph}}{dp}>0,
  \qquad
  \frac{d\zeta_{\rm static}}{dp}\ge0,
  \qquad
  \frac{d\zeta_{\rm NEP}}{dp}<0,
\]
with the degradation attributed to \(\Pi_p(s)\) or \(\Pi_s(s)\), not to the
static graph term alone. That paper should not mix frequency-estimator latency,
stochastic fragility, sheaves, and broad hosting optimization. The minimum
package is a structured Schur derivation, total NEP sensitivity, a
channel-resolved attribution, finite-difference validation, and ablation of
\(L_p\), \(\Pi_p\), and \(\Pi_s\). The empirical headline should be the rate
of wrong-sign static recommendations,
\[
  \Pr\left[
  \operatorname{sign}(\zeta_p^{\rm static})
  \ne
  \operatorname{sign}(\zeta_p^{\rm full})
  \right],
\]
plus the recall achieved by the resolvent-shadow filter when only a small
fraction of line actions are escalated to the full model.

\begin{table}[ht]
\centering
\caption{Minimal result package for the controller-induced reinforcement
reversal paper.}
\label{tab:controller_reversal_paper_package}
\begin{tabular}{L{0.18\textwidth}L{0.34\textwidth}L{0.34\textwidth}}
\toprule
\rowcolor{softblue}
\thead{Result} & \thead{Claim} & \thead{Evidence required}\\
\midrule
A & Exact Schur derivation and mechanism-resolved sensitivity attribution &
Derive \(\Pi(s)\), \(T_s\), and the channel gains \(g_c\).\\
\addlinespace
B & Robust sufficient certificate for controller-induced sign reversal &
Show \(g_0\ge0\) and \(g_0+g_\ell+\bar r<0\) under a registered remainder
bound.\\
\addlinespace
C & Static model wrong-sign rate &
Report strict cases where \(d\zeta_{\rm static}/dp\ge0\) but
\(d\zeta_{\rm full}/dp<0\).\\
\addlinespace
D & Resolvent-shadow computational filter &
Measure precision, recall, false negatives, and cases captured by the top
5--10\% of candidate line actions.\\
\bottomrule
\end{tabular}
\end{table}

\subsection{Passivity Retuning Study}

For a controller parameter \(\kappa\), such as a PLL bandwidth or integral gain,
the core object is
\[
  \mu_\Sigma(\Omega;\kappa)=
  \inf_{\omega\in\Omega}
  \lambda_{\min} H_\Sigma(\omega;\kappa).
\]
The computational question is whether \(\mu_\Sigma\) can be moved from negative
to positive after a stressed operating condition or an \(N-1\) contingency, while
the protected eigenvalue also moves toward the required damping region. The
study is deliberately narrow:
\begin{enumerate}
\item Select a contingency, IBR penetration level, and weak-grid operating
point; re-solve the power flow instead of freezing the pre-contingency
linearization.
\item Sweep \(\kappa\) over the admissible controller range and compute
\(\mu_\Sigma(\Omega;\kappa)\) from the Schur self-energy.
\item Record the active frequency \(\omega_\star\), active Hermitian eigenvector
\(q_\star\), and the envelope derivative
\(q_\star^H(\partial H_\Sigma/\partial\kappa)q_\star\).
\item Identify whether the margin is monotone, locally unimodal, or nonregular
because active frequencies switch.
\item At the best candidate \(\kappa_\star\), recompute the full DAE eigenvalues
and time-domain response. The passivity screen is a gate; the final claim still
needs the full-model check.
\item Compare against a static nodal-\(D\) screen. The negative control is
important: if a scalar damping change predicts the same action, the rational
self-energy did not add design information.
\end{enumerate}

A successful result is not merely a plot where an eigenvalue moves left. It is a
three-part certificate: the self-energy passivity margin improves, the protected
pole moves in the desired direction, and the full model confirms the damping
recovery. Failure is also useful. If \(\mu_\Sigma\) is nonmonotone with multiple
active-frequency switches, the thesis has identified a controller interaction
that cannot be summarized by one PLL retuning knob.

\subsection{Cellular-Sheaf Weak-Link Test}

The sheaf extension has one binary question:
\[
\text{Does a nonconstant sheaf Laplacian reveal a weak link that scalar }
\partial\nu_k/\partial w_e\text{ does not?}
\]
The smallest useful construction is enough for this test. Let each converter bus
carry a stalk \(\mathcal F(v_i)\simeq\R^2\) representing a local PLL phase and
frequency frame. Let each edge stalk \(\mathcal F(e_{ij})\) represent the frame
discrepancy that the line must reconcile. The restriction maps
\(\mathcal F_{v_i\to e_{ij}}\) and \(\mathcal F_{v_j\to e_{ij}}\) are determined
from the solved operating angles and the converter reference frames. With
identity restrictions, this reduces to the constant sheaf and recovers the
ordinary scalar graph intuition. With nontrivial frame rotations, the sheaf
coboundary
\[
  \delta_{\mathcal F}: C^0(G;\mathcal F)\to C^1(G;\mathcal F)
\]
produces the sheaf Laplacian \(L_{\mathcal F}=\delta_{\mathcal F}^\ast
\delta_{\mathcal F}\) on node stalks.

The go--no-go protocol is:
\begin{enumerate}
\item Build \(L_{\mathcal F}\) on the same 39-bus operating point used by the
QEP and Schur studies. The solved angles, not the bare topology, define the
restriction maps.
\item Compare the low spectrum of \(L_{\mathcal F}\) with the scalar
inertia-normalized Laplacian. The comparison should be made under the same
loading and IBR placement assumptions.
\item Remove or down-weight each candidate line and measure the collapse of the
lowest informative sheaf eigenvalues. This gives a sheaf weak-link ranking.
\item Compare that ranking against \(\partial\nu_k/\partial w_e\) and the
existing damping-risk vector. A material ranking difference is the first success
condition.
\item Compute the first harmonic obstruction
\(\dim\ker\delta_{\mathcal F}^\ast\) on high-angle subgraphs. A nonzero value
would indicate a frame-consistency obstruction, i.e., a synchronization
frustration not visible in a scalar Laplacian screen.
\item Validate any sheaf-only weak link in the full model. A sheaf ranking that
does not predict a changed damping, synchronization, or protection-relevant
response is only a mathematical artifact.
\end{enumerate}

The decision rule should be explicit. The sheaf program is a \emph{GO} only if
it produces a weak-link or weak-node ranking that differs materially from the
scalar ranking in at least one sweep case, and the difference has a physical
explanation in frame rotation, PLL alignment, or controller-coordinate
compatibility. It is a \emph{NO-GO} if the sheaf ranking reproduces the scalar
ranking up to scale. In that case the sheaf remains a future mathematical
extension, while the passivity-retuning paper remains the safer near-term
contribution.

\section{Mathematical Audit Gates Before Publication}

The following gates prevent the thesis from overclaiming:
\begin{enumerate}
\item \textbf{Structured Schur gate.} A condensed model must state whether the
rational correction is \(\Pi(s)\), \(\Pi_K(s)\), or \(s\Pi_D(s)\). Calling the
whole Schur correction ``dynamic damping'' is not allowed unless the stiffness
channel has been eliminated or justified.
\item \textbf{Sign-convention gate.} The active-power Jacobian must be tied to
one convention. This thesis uses exported electrical power, so \(L=J_{P\theta}\)
on retained angles.
\item \textbf{Invariant diagnostic gate.} Mode-family labels must use physical
weights, port visibility, left-right participation, or projectors. Raw
controller-state norm ratios are not admissible evidence.
\item \textbf{Curvature gate.} A first-order hosting polytope is a tangent
screen. An inner hosting certificate needs a trust region and a Hessian or model
error bound.
\item \textbf{Cluster gate.} Near repeated or clustered modes, use spectral
projectors and eigenvalue conditioning rather than derivatives of a named mode.
\item \textbf{Resolvent gate.} Controller-pole distance is only a proxy.
Non-normal \(A_{cc}\) requires resolvent gain or port-visible self-energy
dominance.
\item \textbf{Passivity gate.} A Hermitian part of \(\Sigma(j\omega)\) has a
passivity interpretation only after the conjugate ports and supply rate are
declared.
\item \textbf{Total-derivative gate.} A line, inertia, or controller action must
be differentiated through the re-solved operating point whenever the claim is a
planning claim rather than a frozen-linearization screen.
\end{enumerate}

\section{Validation Ladder}

Every claim should climb the same ladder:
\begin{enumerate}
\item \textbf{Algebraic identity.} Verify formulas such as
\(\partial\Lt/\partial M_i\) and the line-shunt decomposition to numerical
precision.
\item \textbf{Surrogate QEP.} Compare diagonal, second-order, reduced-QEP, and
full-QEP damping margins over controlled synthetic systems.
\item \textbf{IEEE topology QEP.} Use realistic network topology and operating
points, while still keeping the reduced second-order abstraction explicit.
\item \textbf{Schur NEP.} Condense controller states and test whether
\(\Pi(s)\), its damping-channel projection \(\Sigma(s)\), and resolvent
dominance explain pole motion.
\item \textbf{Full ANDES eigenanalysis.} Recompute the complete DAE
linearization after each selected action.
\item \textbf{Time-domain and protection gate.} Check RoCoF, nadir,
oscillation decay, estimator exposure, and protection-relevant thresholds.
\end{enumerate}

\begin{figure}[ht]
\centering
\begin{tikzpicture}[
  rung/.style={
    rectangle,
    rounded corners=3pt,
    draw=#1!70!black,
    fill=#1!10,
    text width=0.25\textwidth,
    minimum height=0.78cm,
    align=center,
    font=\small
  },
  arr/.style={-{Latex[length=2.1mm]}, thick, draw=thesisgray}
]
\node[rung=thesisblue] (r1) {1. Algebraic identity};
\node[rung=thesisgreen, right=0.36cm of r1] (r2) {2. Surrogate QEP};
\node[rung=thesisorange, right=0.36cm of r2] (r3) {3. IEEE topology QEP};
\node[rung=thesisred, below=0.55cm of r3] (r4) {4. Schur NEP};
\node[rung=thesisgold, left=0.36cm of r4] (r5) {5. Full ANDES};
\node[rung=thesisblue, left=0.36cm of r5] (r6) {6. Time-domain and protection};
\draw[arr] (r1) -- (r2);
\draw[arr] (r2) -- (r3);
\draw[arr] (r3) -- (r4);
\draw[arr] (r4) -- (r5);
\draw[arr] (r5) -- (r6);
\end{tikzpicture}
\caption{Validation ladder for this research program. A result becomes more
credible as it moves from algebraic identity to full linearization and
time-domain/protection checks.}
\label{fig:validation_ladder_visual}
\end{figure}

Negative gates are part of the thesis. The registered inertia-control cross
term is mathematically real in the local expansion, but the full-ANDES
singleton gate found no measurable effect in 13 tested cases. The correct use
is therefore as a possible second-order diagnostic, not as a headline. The same
discipline applies to hidden margins, line reinforcement, and virtual inertia:
the theory names the mechanism; validation decides whether it is a planning
claim in a given benchmark.
